\section{Introduction}
\label{sec:introduction}

We address the challenge of scheduling business processes in service domains such as finance (e.g., purchasing, order fulfillment), healthcare (e.g., patient flow management, appointment scheduling), and retail (e.g., order processing, supply chain coordination). Business processes are structured sets of activities that organizations execute to achieve various objectives, such as completing a sale, providing a service, or managing a supply chain~\cite{dumas2018fundamentals}. Unlike processes in manufacturing facilities, which typically involve predictable durations, and stable resource availability, business processes exhibit significant uncertainty and variability~\cite{shoush2022intervene,xu2016resource}. 
Specifically, scheduling business processes is notoriously complex due to the stochastic nature of activity durations caused by human behavior, as well as the time-varying availability of resources influenced by factors such as part-time schedules, overlapping responsibilities, and vacations. These challenges make traditional deterministic scheduling methods inadequate~\cite{pinedo2012scheduling}.

In response to this limitation, 
proactive scheduling techniques generate robust offline schedules~\cite{beck2007proactive,chaari2014scheduling}. These methods explicitly account for the stochastic nature of activity durations and aim to compute a \emph{proactive optimal solution} that achieves a pre-defined confidence level~\cite{beck2007proactive}. For example, a solution with a confidence level of \(1-\alpha\) ensures that the returned Makespan remains below an optimal threshold in at least (\(1-\alpha\)) of cases. While proactive scheduling has made significant progress in addressing variability through probabilistic activity durations~\cite{satic2022performance,liu2020multi,hauder2020resource,chen2019research,beck2007proactive}, critical gaps remain. Notably, the challenge of \emph{planned} resource unavailability is seldom addressed. Although studies incorporate the planned unavailability of resources, some overlook the uncertainty of duration~\cite{kreter2018mixed}, while others fail to handle the complexity and variability typical of business processes~\cite{yang2020multi,winklehner2022flexible}.

Beyond methodological gaps, many real-world business processes face an additional practical challenge: the lack of readily available data on key scheduling parameters. Information such as activity durations, their probabilistic distributions, and resource availability calendars is often incomplete or missing. Without access to this data, even advanced scheduling techniques struggle to perform effectively in real-world settings. To bridge this gap, we employ \emph{process mining}, a data-driven approach designed to extract knowledge and insights from event logs~\cite{van2016data}. Event logs are datasets that document process executions, capturing details such as activity types, start and end times, and the resources assigned to each activity. 

Building on this foundation, we formulate the Business Process Scheduling Problem (BPSP) and establish its relationship to the well-known Resource-Constrained Project Scheduling Problems (RCPSP). Our approach leverages both process mining and proactive scheduling methods to tackle both the inherent uncertainty in
business processes, and lack of knowledge of the problem. Specifically, we develop a framework that combines data-driven simulation~\cite{MeneghelloFGR25}, and constraint programming~\cite{kreter2018mixed}: we construct a constraint programming model that generates deterministic schedules for the BPSP, verify them using the simulator, and select the best performing schedules. These schedules are evaluated and optimized to minimize the uncertain Makespan while achieving a pre-specified confidence level, ensuring robustness against variability and resource constraints.

We evaluate our approach in two phases 
to demonstrate both its effectiveness and applicability. In the first phase, we conduct experiments using synthetic data, allowing us to systematically test the framework’s ability to handle varying levels of uncertainty while accounting for resource calendars. This controlled environment ensures that the method can adapt to diverse scheduling scenarios and reliably optimize outcomes under uncertainty. 
In the second phase, we validate the applicability of our approach by applying it to real-world event logs from a healthcare process, achieving an optimization of the process Makespan by an average of 5\% to 14\%.
Using process mining techniques, we extract BPSP parameters from the logs and employ our framework to generate optimal schedules. This highlights the practical value of our method, showcasing its ability to address complex, data-driven scheduling challenges in a real-world domain.


%To summarize, the main contributions of our work are:
%\begin{itemize} 
%\item We formulate the Business Process Scheduling Problem (BPSP) to address the challenges of scheduling processes with stochastic activity durations, and planned resource unavailability. 
%\item We propose a novel BPSP solution framework that combines data-driven business process simulation, leveraging event logs and process mining, with constraint programming to generate robust and optimized schedules under uncertainty. 
%\item We evaluate our framework using synthetic benchmarks and real-world event logs from the healthcare domain. 
%\end{itemize}
